\documentclass[11pt]{article}
\usepackage[a4paper,margin=20mm]{geometry}
\usepackage{amsmath,amssymb,bm,mathtools,physics,microtype,booktabs,enumitem,hyperref}
\setlength{\parindent}{0pt}
\setlength{\parskip}{5pt}
\newcommand{\Id}{\mathbb I}
\newcommand{\Tr}{\operatorname{Tr}}
\newcommand{\ssb}{s\bar s}
\newcommand{\LLb}{\Lambda\bar\Lambda}

\title{\textbf{Daily Research Progress Report -- 22 August 2026}\\
Toward a microscopic dynamical theory of $s\bar s\to\Lambda\bar\Lambda$ spin correlations}
\author{Spin-correlation dynamics project}
\date{22 August 2026}

\begin{document}
\maketitle

\begin{abstract}
The main advance of 22 August is an exact selection rule for the scalar two-spin observable measured in the STAR $\Lambda\bar\Lambda$ analysis. Starting from the most general bilinear two-spin Hamiltonian, we derive the complete coherent equation of motion for the correlation tensor $C_{ij}$. We then prove that common spin precession and any exchange-symmetric bilinear interaction preserve $\Tr C=\langle\bm\sigma_1\cdot\bm\sigma_2\rangle$. Therefore the experimentally observed loss of scalar correlation cannot be explained by such coherent dynamics alone. A reduction of the scalar requires singlet--triplet conversion driven by relative spin fields, exchange-antisymmetric interactions, or genuinely dissipative/collisional terms. This sharply narrows the target of the future two-particle Wigner/Kadanoff--Baym derivation.
\end{abstract}

\section{Physical objective and observable}
We seek a dynamical chain
\begin{equation}
\rho_{s\bar s}^{\rm source}\longrightarrow \rho_{s\bar s}(t)
\longrightarrow \rho_{\Lambda\bar\Lambda}^{\rm obs}
\end{equation}
that connects the initial strange-quark pair to the final hyperon pair. The two-spin density matrix is parameterized as
\begin{equation}
\rho=\frac14\left[\Id\otimes\Id+P_i\sigma_i\otimes\Id+\bar P_j\Id\otimes\sigma_j+C_{ij}\sigma_i\otimes\sigma_j\right].
\end{equation}
The scalar measured after angular averaging is
\begin{equation}
\boxed{P_{\Lambda\bar\Lambda}=\frac13\Tr C=\frac13\langle\bm\sigma_1\cdot\bm\sigma_2\rangle.}
\end{equation}
The STAR short-range central value is approximately
\begin{equation}
P_{\Lambda\bar\Lambda}=0.181\pm0.035_{\rm stat}\pm0.022_{\rm sys},
\end{equation}
with a substantial reduction at larger pair separation.

\section{Initial ${}^3P_0$ source}
For vacuum-like $J^{PC}=0^{++}$ pair creation, the natural state is ${}^3P_0$,
\begin{equation}
L=1,\qquad S=1,\qquad J=0.
\end{equation}
The coupled state is
\begin{equation}
|00\rangle=\frac1{\sqrt3}\sum_{m=-1}^{1}(-1)^{1-m}|1m\rangle_L|1,-m\rangle_S.
\end{equation}
If orbital information is completely traced out,
\begin{align}
\rho_S&=\Tr_L|00\rangle\langle00|\\
&=\frac13\sum_{m=-1}^{1}|1m\rangle\langle1m|\\
&=\frac13P_{S=1}.
\end{align}
Using
\begin{equation}
P_{S=1}=\frac{3\Id+\bm\sigma_1\cdot\bm\sigma_2}{4},
\end{equation}
we obtain
\begin{equation}
\rho_S=\frac14\left[\Id\otimes\Id+\frac13\sum_i\sigma_i\otimes\sigma_i\right],
\end{equation}
so
\begin{equation}
C_{ij}^{(0)}=\frac13\delta_{ij},\qquad P_{s\bar s}^{(0)}=\frac13.
\end{equation}
The distinction between this orbitally averaged state and a momentum-conditioned local state must be retained in a future event-by-event treatment.

\section{Exact coherent evolution of the correlation tensor}
Consider the most general Hermitian bilinear Hamiltonian
\begin{equation}
H=\frac12\Omega_a\sigma_a\otimes\Id+\frac12\bar\Omega_b\Id\otimes\sigma_b+\frac14J_{ab}\sigma_a\otimes\sigma_b.
\end{equation}
The density operator satisfies
\begin{equation}
\dot\rho=-i[H,\rho].
\end{equation}
For
\begin{equation}
C_{ij}=\Tr\left(\rho\,\sigma_i\otimes\sigma_j\right),
\end{equation}
we have
\begin{align}
\dot C_{ij}&=-i\Tr\left([H,\rho]\,\sigma_i\otimes\sigma_j\right)\\
&=-i\left\langle[\sigma_i\otimes\sigma_j,H]\right\rangle,
\end{align}
where cyclicity of the trace was used.

For the first local term,
\begin{align}
-i\left[\sigma_i\otimes\sigma_j,\frac12\Omega_a\sigma_a\otimes\Id\right]
&=-i\frac{\Omega_a}{2}[\sigma_i,\sigma_a]\otimes\sigma_j\\
&=\Omega_a\epsilon_{iak}\sigma_k\otimes\sigma_j,
\end{align}
using $[\sigma_i,\sigma_a]=2i\epsilon_{iak}\sigma_k$. Therefore
\begin{equation}
\left.\dot C_{ij}\right|_1=\epsilon_{iak}\Omega_aC_{kj}.
\end{equation}
Similarly,
\begin{equation}
\left.\dot C_{ij}\right|_2=\epsilon_{jbl}\bar\Omega_bC_{il}.
\end{equation}

For the interaction term we use
\begin{equation}
[A\otimes B,C\otimes D]=\frac12[A,C]\otimes\{B,D\}+\frac12\{A,C\}\otimes[B,D].
\end{equation}
Together with $\{\sigma_i,\sigma_a\}=2\delta_{ia}\Id$ this yields
\begin{equation}
[\sigma_i\otimes\sigma_j,\sigma_a\otimes\sigma_b]
=2i\epsilon_{iak}\delta_{jb}\sigma_k\otimes\Id+2i\delta_{ia}\epsilon_{jbl}\Id\otimes\sigma_l.
\end{equation}
Hence
\begin{equation}
\left.\dot C_{ij}\right|_{12}=\frac12\epsilon_{iak}J_{aj}P_k+\frac12\epsilon_{jbl}J_{ib}\bar P_l.
\end{equation}
Combining all terms gives the exact coherent tensor equation
\begin{equation}
\boxed{\dot C_{ij}=\epsilon_{iak}\Omega_aC_{kj}+\epsilon_{jbl}\bar\Omega_bC_{il}+\frac12\epsilon_{iak}J_{aj}P_k+\frac12\epsilon_{jbl}J_{ib}\bar P_l.}
\end{equation}
No Markov, weak-coupling, or phenomenological channel approximation has entered this derivation.

\section{Protection theorem for the STAR scalar}
Define
\begin{equation}
Q\equiv\bm\sigma_1\cdot\bm\sigma_2=\sum_i\sigma_i\otimes\sigma_i,
\qquad 3P=\langle Q\rangle.
\end{equation}
The spin-swap operator is
\begin{equation}
\mathsf P_{12}=\frac12(\Id+Q),
\end{equation}
so $Q=2\mathsf P_{12}-\Id$. If the Hamiltonian is invariant under spin exchange,
\begin{equation}
\mathsf P_{12}H\mathsf P_{12}=H,
\end{equation}
then
\begin{equation}
[H,\mathsf P_{12}]=0\quad\Rightarrow\quad[H,Q]=0.
\end{equation}
Therefore
\begin{equation}
\boxed{\frac{d}{dt}\langle Q\rangle=0,\qquad \frac{d}{dt}\Tr C=0.}
\end{equation}
For the Hamiltonian above, exchange invariance requires
\begin{equation}
\boxed{\bm\Omega=\bar{\bm\Omega},\qquad J_{ij}=J_{ji}.}
\end{equation}
Thus common precession, isotropic exchange $J\bm S_1\cdot\bm S_2$, and any symmetric anisotropic bilinear coupling may rotate or redistribute tensor components but cannot alter the scalar $P$.

\section{Singlet probability and the microscopic interpretation}
The singlet projector is
\begin{equation}
\Pi_S=|S\rangle\langle S|=\frac14\left(\Id-\bm\sigma_1\cdot\bm\sigma_2\right).
\end{equation}
Hence
\begin{align}
p_S&=\Tr(\rho\Pi_S)\\
&=\frac14\left[1-\langle Q\rangle\right]\\
&=\boxed{\frac{1-3P}{4}}.
\end{align}
A pure triplet source has $p_S(0)=0$ and $P(0)=1/3$. Thus a decrease of the scalar corresponds to increasing singlet weight. Applied only as a hadron-level diagnostic to the STAR central value,
\begin{equation}
p_S^{\rm obs}=\frac{1-3(0.181)}4\simeq0.114.
\end{equation}
This does not yet equal the primordial partonic singlet probability because hadronization and feed-down have not been unfolded.

\section{Coherent mechanisms that can change the scalar}
Write
\begin{equation}
\bm\Omega_c=\frac{\bm\Omega+\bar{\bm\Omega}}2,
\qquad
\delta\bm\Omega=\frac{\bm\Omega-\bar{\bm\Omega}}2.
\end{equation}
Then
\begin{equation}
H_{\rm local}=\bm\Omega_c\cdot(\bm S_1+\bm S_2)+\delta\bm\Omega\cdot(\bm S_1-\bm S_2).
\end{equation}
The first term is exchange symmetric; the second is antisymmetric. For a differential field along $z$,
\begin{equation}
H_\delta=\frac{\delta\Omega_z}{2}(\sigma_{1z}-\sigma_{2z}).
\end{equation}
With
\begin{equation}
|T_0\rangle=\frac{|\uparrow\downarrow\rangle+|\downarrow\uparrow\rangle}{\sqrt2},
\qquad
|S\rangle=\frac{|\uparrow\downarrow\rangle-|\downarrow\uparrow\rangle}{\sqrt2},
\end{equation}
one obtains explicitly
\begin{equation}
(\sigma_{1z}-\sigma_{2z})|T_0\rangle=2|S\rangle.
\end{equation}
Thus $|T_0\rangle$ and $|S\rangle$ mix. Starting from $|T_0\rangle$,
\begin{equation}
p_S(t)=\sin^2(\delta\Omega_z t),
\end{equation}
and therefore
\begin{equation}
\boxed{P(t)=\frac13\left[1-4\sin^2(\delta\Omega_z t)\right].}
\end{equation}

Likewise, the antisymmetric part of $J_{ij}$,
\begin{equation}
J_{ij}^{(A)}=\frac12(J_{ij}-J_{ji})=\epsilon_{ijk}D_k,
\end{equation}
produces
\begin{equation}
H_A=\bm D\cdot(\bm S_1\times\bm S_2),
\end{equation}
which is exchange odd and can also drive singlet--triplet conversion.

\section{Resulting kinetic target}
The scalar problem should therefore be reformulated as singlet--triplet kinetics,
\begin{equation}
\boxed{\dot p_S=\Gamma_{T\to S}p_T-\Gamma_{S\to T}p_S+S_S^{\rm prod},}
\end{equation}
with $p_T=1-p_S$ if no other spin-sector leakage is retained. Since $P=(1-4p_S)/3$,
\begin{equation}
\dot P=-\frac43\left[\Gamma_{T\to S}(1-p_S)-\Gamma_{S\to T}p_S+S_S^{\rm prod}\right].
\end{equation}
The central microscopic task is to derive $\Gamma_{T\to S}$ and $\Gamma_{S\to T}$ from the two-particle QCD collision kernel rather than introduce a generic damping constant.

\section{Connection to relativistic quantum kinetic theory}
A gauge-covariant one-body Wigner function may be decomposed as
\begin{equation}
W=\frac14\left[\mathcal F+i\gamma^5\mathcal P+\gamma^\mu\mathcal V_\mu+\gamma^\mu\gamma^5\mathcal A_\mu+\frac12\sigma^{\mu\nu}\mathcal S_{\mu\nu}\right],
\end{equation}
where $\mathcal A^\mu$ carries spin information. Existing quantum kinetic theory derives spin transport, diffusion, and color-field-induced source terms for massive quarks. For the present problem the required object is the two-particle correlator
\begin{equation}
W^{(2)}(1,2)=\left\langle\widehat W_s(1)\otimes\widehat W_{\bar s}(2)\right\rangle,
\end{equation}
and specifically the connected axial--axial tensor
\begin{equation}
\mathcal C^{\mu\nu}(1,2)=\left\langle\mathcal A_s^\mu(1)\mathcal A_{\bar s}^\nu(2)\right\rangle_c.
\end{equation}
The next derivation should project the two-particle Kadanoff--Baym collision term onto
\begin{equation}
\Pi_S=\frac14(\Id-\bm\sigma_1\cdot\bm\sigma_2),
\qquad
\Pi_T=\Id-\Pi_S,
\end{equation}
with the immediate targets
\begin{equation}
\Gamma_{T\to S}\sim\Tr\left[\Pi_S\,\mathcal C_{\rm coll}[\Pi_T]\right],
\qquad
\Gamma_{S\to T}\sim\Tr\left[\Pi_T\,\mathcal C_{\rm coll}[\Pi_S]\right].
\end{equation}

\section{Hadronization and feed-down}
Even after the partonic kinetics is known, the observable tensor requires a channel-resolved transfer map,
\begin{equation}
C_{ij}^{\Lambda\bar\Lambda}=\sum_{ab}R_{ab}(A_a)_{ik}C_{kl}^{s\bar s,(ab)}(\bar A_b)_{jl}.
\end{equation}
Therefore feed-down should not be represented by a single universal scalar dilution if different hyperon channels have different tensor transfer coefficients.

\section{Conclusions and next steps}
The strongest conclusion of 22 August is not simply that the correlation ``decoheres.'' Instead, the STAR scalar is protected against a broad class of coherent exchange-symmetric dynamics. The relevant microscopic question is which QCD processes create relative spin fields, exchange-antisymmetric couplings, or dissipative collision kernels that transfer probability between triplet and singlet sectors. The next priority is a two-particle Wigner/Kadanoff--Baym derivation of the singlet--triplet transition rates, followed by a source-resolved hadronization/feed-down map into $\Lambda\bar\Lambda$ observables.

\section*{References}
\begin{enumerate}[leftmargin=1.5em]
\item STAR Collaboration, ``Measuring spin correlation between quarks during QCD confinement,'' Nature \textbf{650}, 65--71 (2026).
\item D.-L. Yang, ``Quantum kinetic theory for spin transport of quarks with background chromo-electromagnetic fields,'' JHEP \textbf{06} (2022) 140, arXiv:2112.14392.
\item A. Kumar, B. M\"uller and D.-L. Yang, ``Spin polarization and correlation of quarks from glasma,'' Phys. Rev. D \textbf{107}, 076025 (2023), arXiv:2212.13354.
\item J. Barata, I. Cunt\'in, E. Rico and B. Wu, ``$\Lambda\bar\Lambda$ spin correlations in high-energy collisions from quantum channels: an open quantum system view of hadronization,'' arXiv:2606.30737 (2026).
\end{enumerate}

\end{document}
